\subsection{DIMENSION OF THE HOMOGENEOUS COMPONENTS OF L(X)}

Let X be a set, \(\alpha\) an element of \(\mathbf{N}^{(X)}\) and d an integer >0. We write \(d|\alpha\) if there exists \(\beta\in\mathbf{N}^{(X)}\) such that \(\alpha=d\beta\) . The element \(\beta\) , which is unique, is then denoted by \(\alpha/d\) .

Lemma 1. Let \(n\) be an integer \(>0\) , \(T_1, \ldots, T_n\) indeterminates and \(u_1, \ldots, u_n\) elements of \(\mathbf{Z}\) . Let \((c(\alpha))_{\alpha \in \mathbf{N}^n - \{0\}}\) be a family of elements of \(\mathbf{Z}\) such that

\[
1 - \sum_ {i = 1} ^ {n} u _ {i} T _ {i} = \prod_ {\alpha \neq 0} (1 - T ^ {\alpha}) ^ {c (\alpha)}.\tag{7}
\]

For all \(\alpha \in \mathbf{N}^n - \{0\}\) ,

\[
c (\alpha) = \frac {1}{| \alpha |} \sum_ {d | \alpha} \mu (d) \frac {(| \alpha | / d) !}{(\alpha / d) !} u ^ {\alpha / d}\tag{8}
\]

where \(\mu\) is the Möbius function (Appendix).

Formula (7) is equivalent, on taking logarithms on both sides (Algebra, Chapter IV, § 6, no. 9) to:

\[
\log \left(1 - \sum_ {i = 1} ^ {n} u _ {i} T _ {i}\right) = \sum_ {\alpha \neq 0} c (\alpha) \log (1 - T ^ {\alpha}).\tag{9}
\]

Now

\[
\begin{array}{r l} - \log \left(1 - \sum_ {i = 1} ^ {n} u _ {i} \mathrm{T} _ {i}\right) & = \sum_ {j \geqslant 1} \frac {1}{j} \left(\sum_ {i = 1} ^ {n} u _ {i} \mathrm{T} _ {i}\right) ^ {j} \\ & = \sum_ {j \geqslant 1} \frac {1}{j} \sum_ {| \beta | = j} \frac {| \beta | !}{\beta !} u ^ {\beta} \mathrm{T} ^ {\beta} \\ & = \sum_ {| \beta | > 0} \frac {1}{| \beta |} \frac {| \beta | !}{\beta !} u ^ {\beta} \mathrm{T} ^ {\beta}. \end{array}\tag{10}
\]

On the other hand

\[
\begin{array}{r l} - \sum_ {\alpha \neq 0} c (\alpha) \log (1 - T ^ {\alpha}) & = \sum_ {| \alpha | > 0, k \geqslant 1} \frac {1}{k} c (\alpha) T ^ {k \alpha} \\ & = \sum_ {| \beta | > 0, k | \beta} \frac {1}{k} c \left(\frac {\beta}{k}\right) T ^ {\beta}. \end{array}\tag{11}
\]

Hence (7) is equivalent to

\[
\sum_ {k \mid \beta} \left| \frac {\beta}{k} \right| c \left(\frac {\beta}{k}\right) = \frac {| \beta | !}{\beta !} u ^ {\beta} \quad \text { for   all } \beta \in \mathbf {N} ^ {n} - \{0 \}.\tag{12}
\]

Let \(\Lambda\) be the set of \((\lambda_1, \lambda_2, \ldots, \lambda_n) \in \mathbf{N}^n - \{0\}\) such that the g.c.d. of \(\lambda_1, \lambda_2, \ldots, \lambda_n\) is equal to 1. Every element of \(\mathbf{N}^n - \{0\}\) can be written uniquely

in the form \(m\lambda\) , where \(m\) is an integer \(\geqslant 1\) and \(\lambda \in \Lambda\) . Condition (12) is equivalent to

\[
\sum_ {k \mid m} \left| \frac {m \lambda}{k} \right| c \left(\frac {m \lambda}{k}\right) = \frac {(m | \lambda |) !}{(m \lambda) !} u ^ {m \lambda} \quad \text { for   all } \lambda \in \Lambda \text { and   all } m \geqslant 1.\tag{13}
\]

By the Möbius inversion formula (Appendix), condition (13) is equivalent to

\[
\left| m \lambda \right| c (m \lambda) = \sum_ {d \mid m} \mu (d) \frac {\left| \frac {m \lambda}{d} \right| !}{\left(\frac {m \lambda}{d}\right) !} u ^ {\frac {m \lambda}{d}}\tag{14}
\]

for all \(\lambda\in\Lambda\) and all \(m\geqslant1\) .

THEOREM 2. Let X be a finite set and \(n = \operatorname{Card}(X)\) .

\begin{enumerate}
\def\labelenumi{(\alph{enumi})}
\tightlist
\item
  For every integer \(r \geqslant 1\) , the K-module \(\mathbf{L}^r(\mathbf{X})\) is free of rank
\end{enumerate}

\[
c (r) = \frac {1}{r} \sum_ {d | r} \mu (d) n ^ {r / d},\tag{15}
\]

where μ is the Möbius function.

\begin{enumerate}
\def\labelenumi{(\alph{enumi})}
\setcounter{enumi}{1}
\tightlist
\item
  For all \(\alpha \in \mathbf{N}^{(\mathbf{X})} - \{0\}\) , the K-module \(\mathbf{L}^{\alpha}(\mathbf{X})\) ( \(\S 2\) , no. 6) is free of rank
\end{enumerate}

\[
c (\alpha) = \frac {1}{| \alpha |} \sum_ {d | \alpha} \mu (d) \frac {(| \alpha | / d) !}{(\alpha / d) !}.\tag{16}
\]

We already know that the modules \(L^{r}(\mathbf{X})\) , where \(r \in \mathbf{N}\) , and \(L^{\alpha}(\mathbf{X})\) , where \(\alpha \in \mathbf{N}^{\mathbf{X}}\) , are free (§ 2, no. 11, Corollary to Theorem 1). Consider the multigraduation \((A^{\alpha}(X))_{\alpha \in \mathbf{N}^{\mathbf{X}}}\) of A(X) defined by the canonical homomorphism \(\phi\) of Mo(X) into \(\mathbf{N}^{\mathbf{X}}\) (Algebra, Chapter III, § 3, no. 1, Example 3); then \(A^{\alpha}(\mathbf{X}) \cap L(\mathbf{X}) = L^{\alpha}(\mathbf{X})\) by Remark 3 of no. 1. For \(\alpha \in \mathbf{N}^{\mathbf{X}}\) , the K-module \(A^{\alpha}(\mathbf{X})\) admits as basis the set of words in which each letter \(x\) of X appears \(\alpha(x)\) times. Let \(d(\alpha)\) be the number of these words, that is the rank of \(A^{\alpha}(\mathbf{X})\) ; we shall calculate in two different ways the formal power series

defined by

\[
\mathbf {P} ((\mathrm{T} _ {x}) _ {x \in \mathbf {X}}) \in \mathbf {Z} [ [ (\mathrm{T} _ {x}) _ {x \in \mathbf {X}} ] ]\tag{17}
\]

\[
\mathrm{P} (\mathrm{T}) = \sum_ {\alpha \in \mathbf {N} ^ {\mathbf {X}}} d (\alpha) \mathrm{T} ^ {\alpha}.
\]

\begin{enumerate}
\def\labelenumi{(\arabic{enumi})}
\tightlist
\item
  We have
\end{enumerate}

\[
\mathrm{P} (\mathrm{T}) = \sum_ {m \in \mathrm{Mo} (\mathrm{X})} \mathrm{T} ^ {\phi (m)} = \sum_ {r = 0} ^ {\infty} \sum_ {x _ {1}, \dots , x _ {r}} \mathrm{T} _ {x _ {1}} \dots \mathrm{T} _ {x _ {r}} = \sum_ {r = 0} ^ {\infty} \left(\sum_ {x \in \mathrm{X}} \mathrm{T} _ {x}\right) ^ {r}
\]

whence

\[
\mathrm{P} (\mathrm{T}) = \left(1 - \sum_ {x \in \mathbf {X}} \mathrm{T} _ {x}\right) ^ {- 1}.\tag{18}
\]

\begin{enumerate}
\def\labelenumi{(\arabic{enumi})}
\setcounter{enumi}{1}
\tightlist
\item
  For all \(\alpha \in \mathbf{N}^{\mathbf{X}} - \{0\}\) , let \((e_{\alpha,j})_{1 \leqslant j \leqslant c(\alpha)}\) be a basis of \(\mathrm{L}^{\alpha}(\mathbf{X})\) and give the set I of ordered pairs \((\alpha, j)\) such that \(\alpha \in \mathbf{N}^{\mathbf{X}} - \{0\}\) and \(1 \leqslant j \leqslant c(\alpha)\) a total ordering. By Theorem 1 of no. 1 and the Poincaré-Birkhoff-Witt Theorem (Chapter I, §2, no. 7, Corollary 3 to Theorem 1), the elements
\end{enumerate}

\[
y _ {m} = \prod_ {(\alpha , j) \in \mathbb {I}} (e _ {\alpha , j}) ^ {m (\alpha , j)},
\]

where the index \(\pmb{m}\) runs through \(\mathbf{N}^{(I)}\) , form a basis of \(A(\mathbf{X})\) . Each \(y_{m}\) is of a multidegree \(\sum_{(\alpha, j) \in I} \pmb{m}(\alpha, j)\alpha\) . Let \(u(\pmb{m})\) denote this multidegree. It follows that

\[
\begin{array}{l} \mathrm{P(T)} = \sum_ {m \in \mathrm{N} ^ {(I)}} \mathrm{T} ^ {u (m)} = \sum_ {m \in \mathrm{N} ^ {(I)}} \prod_ {(\alpha , j) \in I} \mathrm{T} ^ {m (\alpha , j) \alpha} \\ = \prod_ {(\alpha , j) \in I} \sum_ {r = 0} ^ {\infty} \mathrm{T} ^ {r \alpha} = \prod_ {(\alpha , j) \in I} (1 - \mathrm{T} ^ {\alpha}) ^ {- 1}, \end{array}
\]

whence finally

\[
\mathrm{P} (\mathrm{T}) = \prod_ {\alpha \in \mathbf {N} ^ {\mathbf {X}} - \{0 \}} (1 - \mathrm{T} ^ {\alpha}) ^ {- c (\alpha)}.\tag{19}
\]

Comparing (18) and (19), we obtain

\[
1 - \sum_ {x \in \mathbf {X}} T _ {x} = \prod_ {\alpha \in \mathbf {N} ^ {\mathbf {X}} - \{0\}} (1 - T ^ {\alpha}) ^ {c (\alpha)}.\tag{20}
\]

Lemma 1 then gives (b).

If we now substitute the same indeterminate U for the \(T_{x}\) for \(x \in X\) in formula (20), we obtain

\[
1 - n \mathrm{U} = \prod_ {\alpha \in \mathbf {N} ^ {\mathrm{X}} - \{0\}} (1 - \mathrm{U} ^ {| \alpha |}) ^ {c (\alpha)} = \prod_ {r > 0} (1 - \mathrm{U} ^ {r}) ^ {c (r)}.
\]

Applying Lemma 1 again, we deduce (a).

Examples. We have

\[
\begin{array}{l l} c (1) = n, & c (2) = \frac {1}{2} (n ^ {2} - n), \qquad c (3) = \frac {1}{3} (n ^ {3} - n), \\ c (4) = \frac {1}{4} (n ^ {4} - n ^ {2}), & c (5) = \frac {1}{5} (n ^ {5} - n), \qquad c (6) = \frac {1}{6} (n ^ {6} - n ^ {3} - n ^ {2} + n). \end{array}
\]

Remark. Let X be a set and let \(\alpha \in \mathbf{N}^{(X)}\) ; the rank of the free K-module \(\mathbf{L}^{\alpha}(X)\) is also given by formula (16). This follows immediately from Theorem 2 (b) and Proposition 4 of § 2, no. 6.

